\documentclass[10pt,a4paper]{amsart}
\usepackage[margin=19mm]{geometry}
\usepackage{amsmath,amssymb,mathtools}
\usepackage[scr=rsfs,cal=euler]{mathalfa}
\usepackage{xfrac}
\usepackage[unicode,hidelinks,bookmarks=false]{hyperref}
\newcommand{\ie}{i.e.\ }
\newcommand{\cf}{cf.\ }
\newcommand{\resp}{respectively\ }
\newcommand{\Subsection}{\subsection}
\providecommand{\nobreakdash}{\nobreak\hbox{-}}
\setlength{\emergencystretch}{2em}
\multlinegap0pt
\usepackage{amssymb}
\usepackage{float}
\usepackage{xcolor}
\newtheorem{remark}{Remark}
\newtheorem{lem}{Lemma}
\newtheorem{prop}{Proposition}
\title{Geometric desingularisation of the sharp-to-smooth travelling wave transition}
\author{Zak Sattar}
\date{Source: arXiv:2608.23489v2 (CC BY 4.0)}
\begin{document}
\maketitle

\noindent\emph{Source and licence.} Geometric desingularisation of the sharp-to-smooth travelling wave transition. Version arXiv:2608.23489v2 (CC BY 4.0), distributed by arXiv under the Creative Commons Attribution 4.0 International licence, http://creativecommons.org/licenses/by/4.0/, as recorded in the arXiv licence metadata for this exact version and shown on the abstract page https://arxiv.org/abs/2608.23489v2. Full licence notice: \texttt{licenses/arxiv-2608.23489-CC-BY-4.0.txt}.\par\smallskip

\noindent\textit{Editorial scope.} Counted unit: the article's prop prop (source0.tex lines 450--452). The article's complete original proof of it is delivered with it (source0.tex lines 453--457, one \texttt{proof} environment of 5 lines). No mathematics is written, repaired or re-derived: the statement and the proof are the article's own lines, in the article's own notation, and no retained line is substituted or re-flowed.  Retained uncounted as the source-local context the proof needs: lines 95--98; lines 144--165; lines 187--206; lines 220--223; lines 248--251; lines 254--257; lines 260--263; lines 283--325; lines 340--364; lines 392--420.  Nothing was omitted from any retained span, so the package carries no omission record.  No retained line contains a citation, so the wrapper carries no bibliography and no bibliography entry.  No retained line contains a figure environment, an image-inclusion command or drawing code, so no image is delivered and the package declares no asset dependency.  Editorial formatting only, in addition: an AMS \texttt{amsart} preamble that restates the article's own declarations and macros used by the retained lines, namely the environments \texttt{remark}, \texttt{lem}, \texttt{prop} on separate counters, together with the standard packages those lines need.  The retained environments print in this document as Remark 1, Lemma 1, Proposition 1 in order of appearance, whereas the article prints the same items with the numbering of its own full document.  The complete native source of the article as distributed in the arXiv e-print for this exact version is bundled unchanged as \texttt{sources/208413/source0.tex}.
\begin{equation}
    u_t=(u^n u_x)_x+u(1-u^n),
    \label{pde1}
\end{equation}

Introducing the travelling wave variable $\eta=x-ct$ and writing $u(x,t)=U(\eta)$, equation \eqref{pde1} becomes
\begin{equation}
    -cU'=(U^nU')'+U(1-U^n).
    \label{tw ode}
\end{equation}
We define $V=U^{n-1}U'$ so that we obtain the system of equations
\begin{equation}
    \begin{aligned}
        U'&=U^{1-n}V,\\
        U^n V'&=-cV-V^2-U^n(1-U^n).
        \label{singular ode}
    \end{aligned}
\end{equation}
We seek solutions which satisfy the boundary conditions $U(-\infty)=1$ and $U(\infty)=0.$ To remove the singularity at $U=0$, we introduce the desingularised variable $\xi$, which is defined through $\tfrac{d}{d\xi}=U^n\tfrac{d}{d\eta}$. Then, system \eqref{singular ode} becomes
\begin{equation}
    \begin{aligned}
        \dot{U} &=UV,\\
        \dot{V} &=-cV-V^2-U^n(1-U^n),
        \label{non-singular ode}
    \end{aligned}
\end{equation}
where the overdot denotes differentiation with respect to the new independent variable $\xi$.

To study smooth front solutions to \eqref{pde1} we perform the preliminary scaling $W=V+\tfrac{1}{\sqrt{1+n}}$, which shifts $\left(0, -c\right)$ to $Q^0\equiv(0, -\varepsilon)$ (or $(0,0)$ in the singular limit), then we set $c=\tfrac{1}{\sqrt{1+n}}+\varepsilon$ and append the trivial $\dot{\varepsilon}=0$, so that \eqref{non-singular ode} becomes
\begin{equation}
    \begin{aligned}
        \dot{U} &=U\left(W-\tfrac{1}{\sqrt{1+n}}\right),\\
        \dot{W} &=\tfrac{1}{\sqrt{1+n}}\varepsilon+\left(\tfrac{1}{\sqrt{1+n}}-\varepsilon\right)W-W^2-U^n(1-U^n),\\
        \dot{\varepsilon}&=0.
        \label{ode epsilon}
    \end{aligned}
\end{equation}
For smooth wavefronts, we require
$$U(-\infty)=1,\;W(-\infty)=\tfrac{1}{\sqrt{1+n}} \;\;\text{and}\;\; U(\infty)=0,\;W(\infty)=\tfrac{1}{\sqrt{1+n}}.$$
\begin{remark}
As $W=V+\tfrac{1}{\sqrt{1+n}}$ and $U^{n-1} U'=V,$
we claim that a smooth front satisfies
$V(\infty)=U^{n-1}(\infty)U'(\infty)=0.$ However, for $n>1$, the condition $V(\infty)=0$ alone does not imply $U'(\infty)=0$, since we could have $U^{n-1}(\infty)=0.$ The required decay rate is given by the
centre-manifold expansion near $(U,W)=\left(0, \tfrac{1}{\sqrt{1+n}}\right)$: $V=W-\tfrac{1}{\sqrt{1+n}}=-\tfrac{1}{\tfrac{1}{\sqrt{1+n}}+\varepsilon}U^n+\mathcal{O}(U^{2n}).$
Consequently,
$U'=-\tfrac{1}{\tfrac{1}{\sqrt{1+n}}+\varepsilon}U+\mathcal{O}(U^{n+1}),$
and therefore,  $U'(\infty)=0$ when $U(\infty)=0$.
\end{remark}

\begin{equation}
    W(U)=\frac{1}{\sqrt{1+n}}U^n+\varepsilon \tfrac{\int_U^1(1-s^n)^{\frac{1+n}{n}}ds}{U(1-U^n)^{\tfrac{1+n}{n}}}+\mathcal{O}(\varepsilon^2).
    \label{perturbed orbit}
\end{equation}

\begin{equation}
    U=r_1, \;\; W=r_1^nw_1,\;\; \varepsilon=r_1^{1+n}\varepsilon_1,
    \label{blow-up in K1}
\end{equation}

\begin{equation}
    U=r_2u_2, \;\; W=r_2^nw_2,\;\; \varepsilon=r_2^{1+n},
    \label{blow-up in K2}
\end{equation}

\begin{equation}
    U=r_3u_3, \;\; W=r_3^n,\;\; \varepsilon=r_3^{1+n}\varepsilon_3.
    \label{blow-up in K3}
\end{equation}

\subsection{$U-$directional chart $K_1$}
In this subsection, we study the dynamics of \eqref{ode epsilon} in the regime $U=\mathcal{O}(1)$. To that end, we consider the $U-$directional chart $K_1$ \eqref{blow-up in K1}. 
Applying \eqref{blow-up in K1} to \eqref{ode epsilon} gives the system 
\begin{equation}
    \begin{aligned}
        \dot{r}_1&=r_1\left(r_1^nw_1-\tfrac{1}{\sqrt{1+n}}\right),\\
        \dot{w}_1&=\tfrac{1}{\sqrt{1+n}}r_1\varepsilon_1+\left(\tfrac{1}{\sqrt{1+n}}-r_1^{1+n}\varepsilon_1\right)w_1-r_1^nw_1^2-(1-r_1^n)-n\left(r_1^nw_1-\tfrac{1}{\sqrt{1+n}}\right)w_1,\\
        \dot{\varepsilon}_1&=-(1+n)\varepsilon_1\left(r_1^nw_1-\tfrac{1}{\sqrt{1+n}}\right).
        \label{K_1}
    \end{aligned}
\end{equation}
Our goal for this subsection is to construct $\Gamma_1$, which is the portion of the singular orbit found in $K_1$.
First notice that \eqref{K_1} admits the line of equilibria  $\ell_1^-=\left\{\left(1, \tfrac{1}{\sqrt{1+n}}, \varepsilon_1\right)\;\bigg|\;\varepsilon_1\in[0, \varepsilon_0)\right\}$. For each fixed $r_1^{1+n}\varepsilon_1=\varepsilon$, $\ell_1^-$ corresponds to the point $Q^-$ before blow-up. As our goal is to construct a singular (in $\varepsilon$) heteroclinic orbit, we identify the point $Q_1^-=\left(1, \tfrac{1}{\sqrt{1+n}}, 0\right)$, obtained by taking $\varepsilon_1\to0$ on $\ell_1^-$. Furthermore, we identify the point $P_1=\left(0,\tfrac{1}{\sqrt{1+n}} ,0\right)$ which corresponds to a projection of $Q^0$ before blow-up. A straightforward calculation gives the following result.
\begin{lem} 
    The equilibrium $Q_1^-$ has the eigenvalues $\tfrac{n}{\sqrt{1+n}}, -\sqrt{1+n}, 0$. The corresponding eigenvectors are $(1,0,0)^T$, $\left(1,-\tfrac{1+2n}{\sqrt{1+n}},0\right)^T$, $(0,0,1)^T$.
    The equilibrium $P_1$ has the eigenvalues $-\tfrac{1}{\sqrt{1+n}}, \sqrt{1+n}, \sqrt{1+n}$. The corresponding eigenvectors are $(1,0,0)^T,(0,1,0)^T, (0,0,1)^T$.
\end{lem}
As $\varepsilon=r_1^{1+n}\varepsilon_1$, there are two limiting systems, namely $\varepsilon_1\to0$ and $r_1\to0$. We first take $\varepsilon_1\to 0$ in \eqref{K_1} to obtain 
\begin{equation}
    \begin{aligned}
        \dot{r}_1&=r_1\left(r_1^n w_1-\tfrac{1}{\sqrt{1+n}}\right),\\
        \dot{w}_1&=\sqrt{1+n}\;w_1-1+r_1^n(1-(1+n)w_1^2).
        \label{K_1 eps_1=0}
    \end{aligned}
\end{equation}
It is obvious that the invariant line $w_1(r_1)=\tfrac{1}{\sqrt{1+n}}$, which we call $\Gamma_1^-$, lies in both $W_1^{\rm u}(Q_1^-)$ and $W_1^{\rm s}(P_1)$. 
The second limiting system is obtained by taking $r_1\to0$ in \eqref{K_1} which yields
\begin{equation}
    \begin{aligned}
        \dot{w}_1&=\sqrt{1+n}w_1-1,\\
        \dot{\varepsilon}_1&=\sqrt{1+n}\varepsilon_1.
        \label{K_1, r_1=0}
    \end{aligned}
\end{equation}
The second portion of the singular orbit in chart $K_1$, which we call $\Gamma_1^+$, must be backward asymptotic to $P_1$ and can be computed by solving 
$\tfrac{d {w}_1}{d{\varepsilon}_1}=\tfrac{\sqrt{1+n}w_1-1}{\sqrt{1+n}\varepsilon_1}$ to obtain $w_1(\varepsilon_1)=\tfrac{1}{\sqrt{1+n}}+\sigma\varepsilon_1$, where $\sigma$ is a constant of integration. To determine $\sigma$, we recall that a perturbation of $W^{\rm u}(Q^-)$ is given by \eqref{perturbed orbit}.  Applying \eqref{blow-up in K1} to \eqref{perturbed orbit}, we find that 
\begin{equation}
    w_1(r_1, \varepsilon_1)=\tfrac{1}{\sqrt{1+n}}+\varepsilon_1
    \frac{\displaystyle \int_{r_1}^1 (1-s^n)^{\tfrac{1+n}{n}}ds}{(1-r_1^n)^{\tfrac{1+n}{n}}}+\mathcal{O}(r_1^{2+n}\varepsilon_1^2), 
    \label{w_1(r_1, eps_1)}
\end{equation}  
then taking $r_1\to0$ gives $w_1(0,\varepsilon_1)=\tfrac{1}{\sqrt{1+n}}+\alpha(n)\varepsilon_1$,
and therefore, $\sigma=\alpha(n)$.

\subsection{Rescaling chart $K_2$}
In this subsection, we study the dynamics of \eqref{ode epsilon} in the regime $U=\mathcal{O}(\sqrt[1+n]{\varepsilon})$ and $W=\mathcal{O}(\sqrt[1+n]{\varepsilon^n})$. To that end, we consider the rescaling chart $K_2$ \eqref{blow-up in K2}.
Applying \eqref{blow-up in K2} to \eqref{ode epsilon} gives 
\begin{equation}
    \begin{aligned}
        \dot{u}_2&=u_2\left(r_2^nw_2-\tfrac{1}{\sqrt{1+n}}\right),\\
        \dot{w}_2&=\tfrac{1}{\sqrt{1+n}}r_2+\left(\tfrac{1}{\sqrt{1+n}}-r_2^{1+n}\right)w_2-r_2^nw_2^2-u_2^n(1-r_2^nu_2^n),\\
        \dot{r}_2&=0.
        \label{K_2}
    \end{aligned}
\end{equation}
Our goal for this subsection is to construct the continuation of $\Gamma_1$ in the chart $K_2$. We shall call this orbit $\Gamma_2$. Notice that \eqref{K_2} admits the line of equilibria $\ell_2^0=(0, -r_2, r_2)$. For each fixed $r_2^{1+n}=\varepsilon$, $\ell_2^0$ corresponds to the point $Q^0$ before blow-up. Furthermore, we identify the point $Q_2^0$ obtained by taking $r_2\to 0$ on $\ell_2^0$. An easy calculation reveals the following.
\begin{lem}
    The origin $Q^0_2=(0,0,0)$ has the eigenvalues $-\tfrac{1}{\sqrt{1+n}}$, $\tfrac{1}{\sqrt{1+n}}, 0$. For $n=1$
    the corresponding eigenvectors are  $\left(1, \tfrac{1}{\sqrt{2}}, 0\right)^T$, $(0,1, 0)^T$, $(0, -1, 1)^T$; for $n>1$, the eigenvectors are $(1,0,0)^T$, $(0,1,0)^T$ and $(0,-1,1)^T$. 
    \label{lem: K_2 linear}
\end{lem}
In this chart, there is only one singular limit to consider, as $\varepsilon=r_2^{1+n}$. Taking $r_2\to0$ in \eqref{K_2} gives 
\begin{equation}
    \begin{aligned}
        \dot{u}_2&=-\tfrac{1}{\sqrt{1+n}}u_2,\\
        \dot{w}_2&=\tfrac{1}{\sqrt{1+n}}w_2-u_2^n.
        \label{K_2 SL}
    \end{aligned}
\end{equation}

\subsection{$W-$directional chart $K_3$}
In this subsection, we study the dynamics of \eqref{ode epsilon} in the regime $W=\mathcal{O}(1)$. To that end, we consider the $W-$directional chart $K_3$ \eqref{blow-up in K3}.
Applying \eqref{blow-up in K3} to \eqref{ode epsilon} gives 
\begin{equation}
    \begin{aligned}
        \dot{u}_3&=u_3\left(r_3^n-\tfrac{1}{\sqrt{1+n}}-\tfrac{1}{n}F(u_3, r_3, \varepsilon_3)\right),\\
        \dot{r}_3&=\tfrac{r_3}{n} F(u_3, r_3, \varepsilon_3),\\
        \dot{\varepsilon}_3&=-\tfrac{1+n}{n}\varepsilon_3 F(u_3, r_3, \varepsilon_3),
        \label{K_3}
    \end{aligned}
\end{equation}
where $F(u_3, r_3, \varepsilon_3)=\tfrac{1}{\sqrt{1+n}}-r_3^n+\tfrac{1}{\sqrt{1+n}}r_3\varepsilon_3-r_3^{1+n}\varepsilon_3-u_3^n(1-r_3^nu_3^n)$.
Our goal for this subsection is to complete the singular orbit, constructing the final part $\Gamma_3$.
Initially, we note that \eqref{K_3} admits the line of equilibria $\ell_3^+=\left\{\left(0, \tfrac{1}{\sqrt[2n]{1+n}}, \varepsilon_3\right)\;\big|\varepsilon_3\in[0, \sqrt[2n]{(1+n)^{1+n}}\;\varepsilon_0)\right\}$. For each fixed $r_3^{1+n}\varepsilon_3=\varepsilon$, $\ell_3^+$ corresponds to the point $Q^+$ before blow-up. Furthermore, we identify the point $Q_3^+=\left(0, \tfrac{1}{\sqrt[2n]{1+n}}, 0\right)$ obtained by taking $\varepsilon_3\to0$ on $\ell_3^+$. Additionally, we note that \eqref{K_3} admits a second relevant equilibrium $P_3=(0,0,0)$, which lies over the degenerate point $(U,W,\varepsilon)=(0,0,0)$ in the blown-up space. An easy calculation reveals the following.
\begin{lem}
    The equilibrium $Q_3^+$ has eigenvalues $0, -\tfrac{1}{\sqrt{1+n}}, 0$. For $n=1$, the corresponding eigenvectors are $(1, -1, 0)^T$, $(0,1,0)^T$ and $(0,0,1)^T$; for $n>1$, the eigenvectors are $(1,0,0)^T, (0,1,0)^T$ and $(0,0,1)^T$.
    The equilibrium $P_3$ has the eigenvalues $-\tfrac{\sqrt{1+n}}{n},\tfrac{1}{n\sqrt{1+n}},-\tfrac{\sqrt{1+n}}{n}$. The corresponding eigenvectors are $(1,0,0)^T, (0,1,0)^T$ and $(0,0,1)^T$.
    \label{lem: K_3 lin}
\end{lem}
As in the chart $K_1$, there are two limiting systems to consider, as $r_3^{1+n}\varepsilon_3=\varepsilon$.
We first take $r_3\to0$, and therefore $F(u_3,0, \varepsilon_3)=\tfrac{1}{\sqrt{1+n}}-u_3^n$, so that \eqref{K_3} becomes
\begin{equation}
    \begin{aligned}
        \dot{u}_3&=\tfrac{u_3}{n}\left(u_3^n-\sqrt{1+n}\right),\\
        \dot{\varepsilon}_3&=-\tfrac{1+n}{n}\varepsilon_3\left(\tfrac{1}{\sqrt{1+n}}-u_3^n\right).
        \label{K_3 r_3=0}
    \end{aligned}
\end{equation}
To solve for $\Gamma_3^-$, we write $\tfrac{d\varepsilon_3}{du_3}=\tfrac{\dot{\varepsilon}_3}{\dot{u}_3}$, giving $\varepsilon_3(u_3)=\gamma u_3\left(1-\tfrac{u_3^n}{\sqrt{1+n}}\right)$.

\begin{prop}
    For \eqref{K_1}, \eqref{K_2} and \eqref{K_3} there exists a singular heteroclinic orbit $\overline{\Gamma}$, connecting $Q_1^-$ and $Q_3^+$.
\end{prop}

\begin{proof}
In $K_1$, the orbit $\Gamma_1^-$ connects $Q_1^-$ to $P_1$, while $\Gamma_1^+$ leaves $P_1$ and intersects $\Sigma_1^{\rm out}$ at $P_1^{\rm out}$. The transition map $\kappa_{12}$ maps this point to
$P_2^{\rm in}$, which determines the orbit $\Gamma_2$ in $K_2$. This orbit intersects $\Sigma_2^{\rm out}$ at $P_2^{\rm out}$. Applying $\kappa_{23}$ gives the point $P_3^{\rm in}$ on $\Gamma_3^-$. The orbit $\Gamma_3^-$ then connects to
$P_3$, and $\Gamma_3^+$ connects $P_3$ to $Q_3^+$. Together these orbits form the singular heteroclinic orbit $\overline{\Gamma}$.
\end{proof}

\end{document}
